\section{Thu nhận media và signaling}

\subsection{Camera và microphone}

Trình duyệt truy cập camera và microphone bằng \texttt{getUserMedia()}. API yêu cầu
HTTPS và quyền của người dùng. Kết quả trả về là một \texttt{MediaStream} chứa các
\texttt{MediaStreamTrack}. Ứng dụng đưa các track này vào
\texttt{RTCPeerConnection} để truyền qua WebRTC \cite{w3c-mediacapture,w3c-webrtc}.

Frontend quản lý quyền thiết bị, thiết bị đang chọn và trạng thái track. Backend không
truy cập camera hoặc microphone. Các trường hợp bắt buộc xử lý gồm từ chối quyền, mất
thiết bị, đổi camera, tắt track và trình duyệt chuyển sang background.

\subsection{RTCPeerConnection và SDP}

\texttt{RTCPeerConnection} quản lý luồng gửi, luồng nhận, ICE transport, codec và trạng
thái kết nối. Hai phía trao đổi offer và answer chứa SDP. SDP mô tả loại media, codec,
hướng truyền, ICE credential và tham số RTP/RTCP \cite{rfc9429}.

Trình tự tối thiểu gồm tạo peer connection, thêm local track, tạo offer, đặt local
description, chuyển offer tới phía còn lại, tạo answer và trao đổi ICE candidate. Remote
track được gắn vào phần tử media khi sự kiện \texttt{track} phát sinh.

\subsection{Vai trò của signaling}

WebRTC không cung cấp kênh signaling. Ứng dụng sử dụng WebSocket, SignalR hoặc HTTPS để
trao đổi offer, answer và ICE candidate \cite{rfc8825}. Kênh này cũng truyền các sự kiện
cuộc gọi: incoming, accept, reject, cancel, busy, timeout và hangup.

Media signaling và call-control signaling phải được phân loại rõ. Backend quản lý trạng
thái nghiệp vụ. Event của trình duyệt phản ánh trạng thái kết nối tại từng endpoint.
